% Part: first-order-logic
% Chapter: tableaux
% Section: identity

\documentclass[../../../include/open-logic-section]{subfiles}

\begin{document}

\olfileid{fol}{tab}{ide}

\olsection{\usetoken{P}{tableau} mawa \usetoken{S}{identity}}

!!^{tableau}s mawa !!{identity} mbutuhake aturan inferensi tambahan.
Aturan kanggo $\eq$ yaiku ($t$, $t_1$, lan $t_2$ iku term-term
tertutup):

\begin{defish}
\AxiomC{}
\RightLabel{$\eq$}
\UnaryInfC{\sFmla{\True}{\eq[t][t]}}
\DisplayProof
\hfill
\AxiomC{\sFmla{\True}{\eq[t_1][t_2]}}
\noLine
\UnaryInfC{\sFmla{\True}{!A(t_1)}}
\RightLabel{$\TRule{\True}{\eq}$}
\UnaryInfC{\sFmla{\True}{!A(t_2)}}
\DisplayProof
\hfill
\AxiomC{\sFmla{\True}{\eq[t_1][t_2]}}
\noLine
\UnaryInfC{\sFmla{\False}{!A(t_1)}}
\RightLabel{$\TRule{\False}{\eq}$}
\UnaryInfC{\sFmla{\False}{!A(t_2)}}
\DisplayProof
\end{defish}
Gatekna yen, beda karo kabeh aturan liyane, $\TRule{\True}{\eq}$ lan
$\TRule{\False}{\eq}$ mbutuhake \emph{rong} !!{formula}s wis
muncul ing cabang, yaiku $\sFmla{\True}{\eq[t_1][t_2]}$ lan
$\sFmla{S}{!A(t_1)}$ kalorone.

\begin{ex}
Yen $s$ lan $t$ iku term-term tertutup, mula $\eq[s][t], !A(s)
\Proves !A(t)$:
\begin{oltableau}
  [\sFmla{\False}{\formula{A}(t)}, just = \TAss
    [\sFmla{\True}{\eq[s][t]}, just = \TAss
      [\sFmla{\True}{\formula{A}(s)}, just = \TAss
        [\sFmla{\True}{\formula{A}(t)}, just={\TRule{\True}{\eq}[2, 3]}, close]
      ]
    ]
  ]
\end{oltableau}
Iki bisa wae wis dikenal minangka prinsip substitutabilitas barang
kang identik, utawa Hukum Leibniz.

!!^{tableau}s mbuktekake yen $\eq$ iku simetris, yaiku
$\eq[s_1][s_2] \Proves \eq[s_2][s_1]$:
\begin{oltableau}
  [\sFmla{\False}{\eq[s_2][s_1]}, just = \TAss
    [\sFmla{\True}{\eq[s_1][s_2]}, just = \TAss
      [\sFmla{\True}{\eq[s_1][s_1]}, just = {$\eq$}
        [\sFmla{\True}{\eq[s_2][s_1]}, just = {\TRule{\True}{\eq}[2, 3]}, close]
      ]
    ]
  ]
\end{oltableau}
Ing kene, larik~$2$ iku !!{formula} prasyarat kapisan
$\sFmla{\True}{\eq[s_1][s_2]}$ saka $\TRule{\True}{\eq}$.
Larik~$3$ iku prasyarat kapindho, awujud
$\sFmla{\True}{!A(s_1)}$---anggepen $!A(x)$ minangka
$\eq[x][s_1]$; mula $!A(s_1)$ yaiku $\eq[s_1][s_1]$ lan $!A(s_2)$
yaiku $\eq[s_2][s_1]$. Cathetan koreksi sumber OLPL-021: sumber
Inggris pisanan menehi argumen kapindho marang formula prasyarat ing
larik katelu, nanging substitusi kang langsung diterangake sawise iku
lan arah aturan netepake argumen kapisan, kaya kang dibalekake ing
kene.

Tableaux uga mbuktekake yen $\eq$ iku transitif, yaiku
$\eq[s_1][s_2], \eq[s_2][s_3] \Proves \eq[s_1][s_3]$:
\begin{oltableau}
  [\sFmla{\False}{\eq[s_1][s_3]}, just = \TAss
    [\sFmla{\True}{\eq[s_1][s_2]}, just = \TAss
      [\sFmla{\True}{\eq[s_2][s_3]}, just = \TAss
        [\sFmla{\True}{\eq[s_1][s_3]}, just = {\TRule{\True}{\eq}[3, 2]}, close]
      ]
    ]
  ]
\end{oltableau}
Ing !!{tableau} iki, !!{formula} prasyarat kapisan saka
$\TRule{\True}{\eq}$ yaiku larik~$3$,
$\sFmla{\True}{\eq[s_2][s_3]}$ ($s_2$ nindakake peran~$t_1$, lan
$s_3$ peran~$t_2$). Prasyarat kapindho, kang awujud
$\sFmla{\True}{!A(s_2)}$, yaiku larik~$2$. Ing kene, anggepen
$!A(x)$ minangka $\eq[s_1][x]$; iki nggawe $!A(s_2)$ dadi
$\eq[s_1][s_2]$ (yaiku larik~$2$) lan $!A(s_3)$ dadi
!!{formula}~$\eq[s_1][s_3]$ ing kesimpulan. Cathetan koreksi sumber
OLPL-022: sumber Inggris malih menehi pasangan metavariabel aturan
kanggo larik kapindho, sanadyan penetapan kang wis kasebut bakal
menehi larik katelu; formula larik kapindho kang dibalekake ing kene
ditemtokake dening tableau lan pilihan formula substitusine.
\end{ex}

\begin{prob}
Wenehana !!{tableau}s katutup kanggo perkara-perkara iki:
\begin{enumerate}
\item $\sFmla{\False}{\lforall[x][\lforall[y][((x = y \land !A(x))
      \lif !A(y))]]}$
\item $\sFmla{\False}{\lexists[x][(!A(x) \land
    \lforall[y][(!A(y) \lif y = x)])]}$,\\
  $\sFmla{\True}{\lexists[x][!A(x)] \land
    \lforall[y][\lforall[z][((!A(y) \land !A(z)) \lif y = z)]]}$
\end{enumerate}
\end{prob}

\end{document}
